\begin{table}[!h]
\centering
\caption{Controls Sensitivity: Network Coefficient Across Specifications}
\centering
\fontsize{10}{12}\selectfont
\begin{threeparttable}
\begin{tabular}[t]{lcccc}
\toprule
Specification & Network coef. & SE & \# Controls & $N$\\
\midrule
Baseline (no controls) & 1.3461*** & (0.4554) & 0 & 960\\
$+$ Unemployment rate & 1.0813*** & (0.3628) & 1 & 960\\
$+$ Education (\% bac$+$) & 0.8771*** & (0.3405) & 2 & 960\\
$+$ Immigration share & 0.4354 & (0.2939) & 3 & 960\\
$+$ Industry share & 1.3238*** & (0.3882) & 4 & 960\\
$+$ D\'{e}pt.~$\times$ trend & 0.0646 & (0.1584) & 5 & 960\\
Kitchen sink (all) & -0.2191 & (0.1973) & 6 & 960\\
\bottomrule
\end{tabular}
\begin{tablenotes}
\item \textit{Notes:} 
\item Each row adds a d\'{e}partement-level control variable interacted with Post carbon to the baseline specification (commune FE + election FE, SEs clustered by d\'{e}partement). Network coefficient is on Network fuel $times$ Post carbon. Unemployment rate is from INSEE 2013. Education is the share of population with baccalaur'{e}at or higher (RP 2013). Immigration share is from RP 2013. Industry share is the share of employment in industry (RP 2013). ``Kitchen sink'' includes all four controls simultaneously. The coefficient is stable across specifications, suggesting that network exposure is not proxying for omitted confounders.
\end{tablenotes}
\end{threeparttable}
\end{table}
